\section{提案手法}\label{sec:method}

本節では、密集区間予測のための3段階パイプラインを提案する：
(1) IDIT分布分析、(2) Hawkes-Denseモデルによる発生時刻予測、
(3) Weibull生存モデルによる持続時間予測。

\subsection{IDIT分布分析}

密集区間間時間 $\{\tau_1, \tau_2, \ldots, \tau_{K-1}\}$ に対して、
以下の4つの候補分布をフィットし、
AIC~\cite{akaike1974new}およびBIC~\cite{schwarz1978estimating}で選択する。

\begin{itemize}
  \item \textbf{指数分布}（Exponential）：$f(\tau) = \lambda e^{-\lambda \tau}$。
    密集区間の発生がPoisson過程に従う場合に適合する。
  \item \textbf{Weibull分布}：$f(\tau) = \frac{k}{\lambda}\left(\frac{\tau}{\lambda}\right)^{k-1} e^{-(\tau/\lambda)^k}$。
    形状パラメータ $k$ により経時変化を表現できる。
  \item \textbf{ガンマ分布}：$f(\tau) = \frac{\beta^\alpha}{\Gamma(\alpha)} \tau^{\alpha-1} e^{-\beta \tau}$。
    柔軟な形状を持つ。
  \item \textbf{対数正規分布}：$f(\tau) = \frac{1}{\tau \sigma \sqrt{2\pi}} e^{-\frac{(\ln \tau - \mu)^2}{2\sigma^2}}$。
    右裾が重い分布に適合する。
\end{itemize}

適合度はKolmogorov-Smirnov検定により評価する。

\subsection{Hawkes-Denseモデル}

密集区間の到着時刻列 $\{s_1, s_2, \ldots, s_K\}$ を
自己励起点過程でモデル化する。
条件付き強度関数を
\begin{equation}\label{eq:hawkes-intensity}
  \lambda(t|\mathcal{H}_t) = \mu + \sum_{s_k < t} \alpha \beta \exp(-\beta(t - s_k))
\end{equation}
と定義する。
ここで $\mu > 0$ はベースライン強度、$0 < \alpha < 1$ は自己励起パラメータ
（分岐比; branching ratio）、$\beta > 0$ は減衰率である。

$\alpha < 1$ の条件下で過程は定常であり、
定常強度は $\bar{\lambda} = \mu / (1 - \alpha)$ となる。

\subsubsection{パラメータ推定}

対数尤度関数は点過程の一般理論~\cite{daley2003introduction}より
\begin{equation}\label{eq:loglikelihood}
  \ell(\mu, \alpha, \beta) = \sum_{k=1}^{K} \log \lambda(s_k|\mathcal{H}_{s_k}) - \int_0^T \lambda(t|\mathcal{H}_t) \, dt
\end{equation}
で与えられる。補償項（integral）は解析的に計算可能で、
\begin{equation}
  \int_0^T \lambda(t|\mathcal{H}_t) \, dt = \mu T + \alpha \sum_{k=1}^{K} \left(1 - e^{-\beta(T - s_k)}\right)
\end{equation}
となる。Nelder-Mead法により最尤推定を行う。

\subsubsection{次回発生時刻予測}

推定されたパラメータを用いて、
Ogataのthinningアルゴリズム~\cite{ogata1981lewis}により
次回の密集区間発生時刻をシミュレーションする。

\begin{proposition}[予測の一致性]\label{prop:consistency}
  以下の条件を仮定する：
  (i) 真のパラメータ $(\mu^*, \alpha^*, \beta^*)$ は $\mu^* > 0$, $0 \leq \alpha^* < 1$, $\beta^* > 0$ を満たす、
  (ii) 観測区間 $[0, T]$ において $T \to \infty$ かつ $K/T \to \bar{\lambda}$（定常強度）。
  このとき、最尤推定量 $(\hat{\mu}, \hat{\alpha}, \hat{\beta})$ は
  真のパラメータ $(\mu^*, \alpha^*, \beta^*)$ に確率収束する。
  さらに、thinningアルゴリズムによる予測分布
  $\hat{F}_{K+1}(t) = P(s_{K+1} \leq t | \hat{\mu}, \hat{\alpha}, \hat{\beta}, \mathcal{H}_{s_K})$ は
  真の条件付き分布 $F^*_{K+1}(t)$ に弱収束する。
\end{proposition}

\begin{proof}[証明のスケッチ]
  (1) 定常Hawkes過程の最尤推定量の一致性と漸近正規性は
  Ogata~\cite{ogata1978asymptotic}により示されている。
  具体的には、$\sqrt{K}(\hat{\theta} - \theta^*) \xrightarrow{d} \mathcal{N}(0, I(\theta^*)^{-1})$
  が成り立つ（$I(\theta^*)$ はFisher情報行列）。
  (2) thinningアルゴリズムは条件付き強度関数 $\lambda(t|\mathcal{H}_t)$ からの
  正確なシミュレーション手法であり~\cite{ogata1981lewis}、
  $\lambda$ の連続性から
  $\hat{\theta} \xrightarrow{p} \theta^*$ は
  $\hat{F}_{K+1} \Rightarrow F^*_{K+1}$ を導く。

  \textbf{注意}：$\alpha^* = 0$ の場合、Hawkes過程はPoisson過程に退化し、
  $\alpha$ の最尤推定量はパラメータ空間の境界に位置するため、
  通常の漸近理論は直接適用されない。
  この場合、尤度比検定で自己励起構造の有意性を検定する必要がある（\S\ref{sec:experiments}参照）。
\end{proof}

\subsection{持続時間予測}

密集区間の持続時間 $\{d_1, d_2, \ldots, d_K\}$ を
Weibull分布~\cite{weibull1951statistical}でモデル化する。

生存関数は
\begin{equation}
  S(d) = \exp\left(-\left(\frac{d}{\lambda}\right)^k\right)
\end{equation}
で与えられ、ハザード関数は
\begin{equation}
  h(d) = \frac{k}{\lambda}\left(\frac{d}{\lambda}\right)^{k-1}
\end{equation}
となる。

形状パラメータ $k$ の値により以下の解釈が得られる：
\begin{itemize}
  \item $k < 1$：時間とともにハザードが減少（長い密集区間ほど終了しにくい）
  \item $k = 1$：指数分布（無記憶性; 終了確率は持続時間に依存しない）
  \item $k > 1$：時間とともにハザードが増加（長い密集区間ほど終了しやすい）
\end{itemize}

予測持続時間は、フィットされた分布の平均
$\hat{d} = \lambda \Gamma(1 + 1/k)$
または中央値
$\tilde{d} = \lambda (\ln 2)^{1/k}$
として与えられる。
信頼区間は分位点から計算する。

\subsection{統合パイプライン}

提案手法の統合パイプラインをAlgorithm~\ref{alg:pipeline}に示す。

\begin{algorithm}[t]
\caption{Dense Interval Prediction Pipeline}
\label{alg:pipeline}
\begin{algorithmic}[1]
\Require トランザクションDB $\mathcal{D}$, ウィンドウサイズ $W$, 閾値 $\theta$
\Ensure 次回発生時刻予測 $\hat{s}_{K+1}$, 持続時間予測 $\hat{d}_{K+1}$
\State $\mathcal{F} \gets \textsc{DenseItemsetMining}(\mathcal{D}, W, \theta)$
\For{each $(X, \mathcal{I}_X) \in \mathcal{F}$}
  \State $\{\tau_k\} \gets \textsc{ComputeIDIT}(\mathcal{I}_X, W)$
  \State $F^* \gets \textsc{FitIDIT}(\{\tau_k\})$ \Comment{AIC/BICで最良分布選択}
  \State $(\hat{\mu}, \hat{\alpha}, \hat{\beta}) \gets \textsc{FitHawkes}(\{s_k\})$
  \State $\hat{s}_{K+1} \gets \textsc{ThinningPredict}(\hat{\mu}, \hat{\alpha}, \hat{\beta}, \{s_k\})$
  \State $(\hat{k}, \hat{\lambda}) \gets \textsc{FitWeibull}(\{d_k\})$
  \State $\hat{d}_{K+1} \gets \hat{\lambda} \Gamma(1 + 1/\hat{k})$
\EndFor
\end{algorithmic}
\end{algorithm}

\begin{theorem}[計算量]
  提案パイプラインの計算量は、$K$ 個の密集区間に対して、
  IDIT分布フィッティングが $O(K \log K)$、
  Hawkes過程のパラメータ推定が $O(K^2 I)$（$I$ は最適化反復回数）、
  thinning予測が $O(KM)$（$M$ はサンプル数）、
  Weibullフィッティングが $O(K)$ である。
  したがって全体の計算量は $O(K^2 I + KM)$ である。
\end{theorem}
